\begin{lemma}
\label{lemma-lift-section-smooth-morphism}
Let $A$ be a ring, let $I \subset A$ be an ideal. Consider a commutative
diagram
$$
\xymatrix{
B \ar[rd] \\
A \ar[u] \ar[r] & A/I
}
$$
where $B$ is a smooth $A$-algebra. Then there exists an \'etale ring
map $A \to A'$ which induces an isomorphism $A/I \to A'/IA'$ and an
$A$-algebra map $B \to A'$ lifting the ring map $B \to A/I$.
\end{lemma}

\begin{proof}
Let $J \subset B$ be the kernel of $B \to A/I$ so that $B/J = A/I$. By
Algebra, Lemma \ref{algebra-lemma-application-NL-smooth} the sequence
$$
0 \to I/I^2 \to J/J^2 \to \Omega_{B/A} \otimes_B B/J \to 0
$$
is split exact. Thus $\overline{P} = J/(J^2 + IB) = \Omega_{B/A} \otimes_B B/J$
is a finite projective $A/I$-module. Choose an integer $n$ and a direct sum
decomposition $A/I^{\oplus n} = \overline{P} \oplus \overline{K}$.
By Lemma \ref{lemma-lift-projective-module} we can find an
\'etale ring map $A \to A'$ which induces an isomorphism
$A/I \to A'/IA'$ and a finite projective $A$-module $K$ which
lifts $\overline{K}$. We may and do replace $A$ by $A'$.
Set $B' = B \otimes_A \text{Sym}_A^*(K)$. Since $A \to \text{Sym}_A^*(K)$
is smooth by Lemma \ref{lemma-symmetric-algebra-smooth} we see that
$B \to B'$ is smooth which in turn implies that $A \to B'$ is smooth (see
Algebra, Lemmas \ref{algebra-lemma-base-change-smooth} and
\ref{algebra-lemma-locally-smooth}).
Moreover the section $\text{Sym}^*_A(K) \to A$ determines a section
$B' \to B$ and we let $B' \to A/I$ be the composition $B' \to B \to A/I$.
Let $J' \subset B'$ be the kernel of $B' \to A/I$. We have
$JB' \subset J'$ and $B \otimes_A K \subset J'$. These maps combine
to give an isomorphism
$$
(A/I)^{\oplus n} \cong J/(J^2 + IB) \oplus \overline{K}
\longrightarrow
J'/((J')^2 + IB')
$$
Thus, after replacing $B$ by $B'$ we may assume that
$J/(J^2 + IB) = \Omega_{B/A} \otimes_B B/J$ is a free
$A/I$-module of rank $n$.

\medskip\noindent
In this case, choose $f_1, \ldots, f_n \in J$ which map to a
basis of $J/(J^2 + IB)$. Consider the finitely presented $A$-algebra
$C = B/(f_1, \ldots, f_n)$. Note that we have an exact sequence
$$
0 \to H_1(L_{C/A}) \to (f_1, \ldots, f_n)/(f_1, \ldots, f_n)^2
\to \Omega_{B/A} \otimes_B C \to \Omega_{C/A} \to 0
$$
see Algebra, Lemma \ref{algebra-lemma-exact-sequence-NL} (note that
$H_1(L_{B/A}) = 0$ and that $\Omega_{B/A}$ is finite projective,
in particular flat so the Tor group vanishes). For any prime
$\mathfrak q \supset J$ of $B$ the module $\Omega_{B/A, \mathfrak q}$
is free of rank $n$ because $\Omega_{B/A}$ is finite projective and
because $\Omega_{B/A} \otimes_B B/J$ is free of rank $n$
(see Algebra, Lemma \ref{algebra-lemma-finite-projective}). By our choice
of $f_1, \ldots, f_n$ the map
$$
\left((f_1, \ldots, f_n)/(f_1, \ldots, f_n)^2\right)_{\mathfrak q}
\to
\Omega_{B/A, \mathfrak q}
$$
is surjective modulo $J$. Hence we see that this map of modules over
the local ring $C_{\mathfrak q}$ has to be an isomorphism
(by Algebra, Lemma \ref{algebra-lemma-NAK} the map is surjective
and for example by Algebra, Lemma \ref{algebra-lemma-fun}
because $((f_1, \ldots, f_n)/(f_1, \ldots, f_n)^2)_{\mathfrak q}$
is generated by $n$ elements the map is injective). Thus
$H_1(L_{C/A})_{\mathfrak q} = 0$ and $\Omega_{C/A, \mathfrak q} = 0$. By
Algebra, Lemma \ref{algebra-lemma-smooth-at-point}
we see that $A \to C$ is smooth at the prime $\overline{\mathfrak q}$
of $C$ corresponding to $\mathfrak q$. Since
$\Omega_{C/A, \mathfrak q} = 0$ it is actually \'etale at
$\overline{\mathfrak q}$. Thus $A \to C$ is \'etale at all primes
of $C$ containing $JC$. By Lemma \ref{lemma-localize-upstairs}
we can find an $f \in C$ mapping to an invertible element of $C/JC$ such that
$A \to C_f$ is \'etale. By our choice of $f$ it is still true that
$C_f/JC_f = A/I$. The map $C_f/IC_f \to A/I$ is surjective and
\'etale by Algebra, Lemma \ref{algebra-lemma-map-between-etale}.
Hence $A/I$ is isomorphic to the localization of $C_f/IC_f$ at
some element $g \in C$, see
Algebra, Lemma \ref{algebra-lemma-surjective-flat-finitely-presented}.
Set $A' = C_{fg}$ to conclude the proof.
\end{proof}